\pagebreak

\section{Xilinx VirtexII Pro MGT library - xilinx/mgt.3pl}

    \index{library!xilinx/mgt.3pl}
    This library provide a simple 16 bit interface to a VirtexII Pro
    Multi-Gigabit Transceiver (referred to as an MGT or Rocket).
    
    This library must be incorporated by using a module statement, e.g. -
    \begin{lstlisting}[gobble=4, language=threepl]
       module "xilinx/mgt.3pl"
    \end{lstlisting}
    
    The global structures {\bf mgt\_16}, {\bf mgt\_16\_e} and {\bf mgt\_16\_ee}
    are created. They are defined as follows -

    \begin{lstlisting}[gobble=4, language=threepl]
       struct(global.gt_16,
           data,   "bits:16",  // data
           k,      "log",      // true if top 8 bits of data form a K-character
           contin, "log"       // this is a continuation word
       );
    \end{lstlisting}

    \begin{lstlisting}[gobble=4, language=threepl]
       struct(global.gt_16_e,
           data,   "bits:16",  // data
           k,      "log",      // true if top 8 bits of data form a K-character
           first,  "log",      // previous word was an idle
           error,  "log"       // disparity error or not-in-table error
       );
    \end{lstlisting}

    \begin{lstlisting}[gobble=4, language=threepl]
       struct(global.gt_16_ee,
           data,       "bits:16",  // data
           k,          "log",      // true if top 8 bits of data form a K-character
           first,      "log",      // previous word was an idle
           disperr,    "[2]log",   // disparity errors
           niterr,     "[2]log"    // not-in-table errors
       );
    \end{lstlisting}

    {\bf mgt\_16} is used for input data. {\bf mgt\_16\_e} and
    {\bf mgt\_16\_ee} are used for output data, the former providing a simpler
    error indication than the latter.
    
    The global structure {\bf ds16} is also used and it is created in
    {\bf stdlib.3pl}. It is 
    \begin{lstlisting}[gobble=4, language=threepl]
       struct(global.ds16,
           data,   "bits:16",  // data
           k,      "log"       // true if top 8 bits of data form a K-character
       );
    \end{lstlisting}


    \subsection{library modules}

        \subsubsection{gt\_custom\_16 - a 16-bit Rocket high-speed serial I/O channel}\label{sec:lmgtcustom16}

            {\bf Library: xilinx/mgt.3pl}

            {\bf Prototype:} \\
            \vsp
            \begin{lstlisting}[gobble=12, language=threepl]
               global module
                   gt_custom_16 (
                       queue(out),
                       clock(ruclk),
                       value(rxdiag, "[5]log"),
                       value(txdiag, "[2]log")
                       <-
                       str({loc, txp_pin, txn_pin, rxp_pin, rxn_pin, clkp_pin, clkn_pin}),
                       float(freq),         // MGT base frequency
                       queue(in),
                       int(run_length),
                       value(reset, "log"),
                       int(tx_diff_ctrl, 500),
                       int(tx_preemphasis, 0),
                       value(loopback, "uint:2", 0),
                       log(los_fsm)
                   ) {}
            \end{lstlisting}

            {\bf Output parameters:} \\
            {\bf out} is the data output.\\
            {\bf ruclk} is the user interface clock.\\
            {\bf rxdiag} provides five diagnostic receiver signals.\\
            {\bf txdiag} provides two diagnostic transmitter signals.

            {\bf Input parameters:} \\
            {\bf loc} gives the MGT location coordinates.\\
            {\bf txp\_pin} specifies the transmitter positive output pin.\\
            {\bf txn\_pin} specifies the transmitter negative output pin.\\
            {\bf rxp\_pin} specifies the receiver positive input pin.\\
            {\bf rxn\_pin} specifies the receiver negative input pin.\\
            {\bf clkp\_pin} specifies the clock input positive pin.\\
            {\bf clkn\_pin} specifies the clock input negative pin.\\
            {\bf freq} is the parallel input and output data frequency.\\
            {\bf in} is the data input.\\
            {\bf run\_length} is the maximum data run length.\\
            {\bf reset} when true resets the initialisation flag.\\
            {\bf tx\_diff\_ctrl} specifies the differential output voltage.\\
            {\bf tx\_preemphasis} specifies the transmitter preemphasis.\\
            {\bf loopback} when true connects the MGT in loopback mode.\\
            {\bf los\_fsm} when true connects enables the MGT loss-of-sync FSM.

            {\bf Description:} \\
            This module provides a simple interface to procedure
            {\bf mgt()} \autorefph{sec:lpmgt}, in this library.
            It creates both an input and an output interface to the
            VirtexII Pro MGT. The data width is 16 bit. A custom protocol
            is used. 
            
            The particular MGT is selected by the first 7 input arguments.
            The {\bf loc} does not seem to be required but the transmitter pin,
            receiever pin and clock pin pairs must be provided. The following
            combinations apply to the XC2VP4 -           
            
            \begin{tabular}{| l || r | r | r | r |}
            \hline
            MGT      & TXP  & TXN  & RXP  & RXN  \\
            location & pin  & pin  & pin  & pin  \\
            \hline
            \hline
            GT\_X0Y0 & AB8  & AB7  & AB9  & AB10 \\
            \hline
            GT\_X0Y1 & A8   & A7   & A9   & A10  \\
            \hline
            GT\_X1Y0 & AB14 & AB13 & AB15 & AB16 \\
            \hline
            GT\_X1Y1 & A14  & A13  & A15  & A16  \\
            \hline
            \end{tabular}
            
            and for the XC2VP40 -

            \begin{tabular}{| l || r | r | r | r |}
            \hline
            MGT      & TXP  & TXN  & RXP  & RXN  \\
            location & pin  & pin  & pin  & pin  \\
            \hline
            \hline
            GT\_X0Y0 & ap32 & ap33 & ap31 & ap30 \\
            \hline
            GT\_X0Y1 &  a32 &  a33 &  a31 &  a30 \\
            \hline
            GT\_X1Y0 & ap28 & ap29 & ap27 & ap26 \\
            \hline
            GT\_X1Y1 &  a28 &  a29 &  a27 &  a26 \\
            \hline
            GT\_X2Y0 & ap20 & ap21 & ap19 & ap18 \\
            \hline
            GT\_X2Y1 &  a20 &  a21 &  a19 &  a18 \\
            \hline
            GT\_X3Y0 & ap16 & ap17 & ap15 & ap14 \\
            \hline
            GT\_X3Y1 &  a16 &  a17 &  a15 &  a14 \\
            \hline
            GT\_X4Y0 &  ap8 &  ap9 &  ap7 &  ap6 \\
            \hline
            GT\_X4Y1 &   a8 &   a9 &   a7 &   a6 \\
            \hline
            GT\_X5Y0 &  ap4 &  ap5 &  ap3 &  ap2 \\
            \hline
            GT\_X5Y1 &   a4 &   a5 &   a3 &   a2 \\
            \hline
            \end{tabular}
            
            {\bf freq} is an immediate floating point value which specifies the
            MGT base frequency in Hz. For 2.5GHz serial the base frequency
            is 125.0MHz (serial rate divided by 20).
            
            Input parameter {\bf in} is the input data queue for the transmitter.
            If the transmitter is not used the argument is omitted. The type is
            either the struct {\bf ds16} (defined in stdlib.3pl) or the struct
            {\bf gt\_16}. The most significant byte of the word may be a
            K-character, which is signified by setting field {\bf k} true. The
            byte must be one of the allowed K-characters (this is not checked!).
            For type {\bf gt\_16} the {\bf contin} field should be true if this
            word is a continuation word, i.e. it is not the first word in a
            block. This ensures that despite the run length constraint (see
            Input parameter {\bf run\_length} below) a block will not be split
            by inserted idles. At the receiving end the first word following an
            idle can reset a block boundary.

            Input parameter {\bf run\_length} is an immediate integer which
            specifies the maximum contiguous data run for the transmitter. If
            data input is continuously available such that this length would be
            exceeded (and the data {\bf contin} field is not true - see above)
            then an idle sequence is inserted in the output stream. When no data
            are available for transmission a continuous stream of idles is
            transmitted. These idles are used by the receiver to lock initially
            and to continue checking that lock is maintained. If receiver lock
            is lost then lock can only be regained at the next idle. In the
            custom protocol used in this interface the idle sequence is a single
            16-bit word containing a comma byte. If no argument is supplied
            to parameter {\bf run\_length} then idles will only be inserted
            when the input queue is unavailable.

            Input parameter {\bf reset} is used to clear the initialisation flag
            and execute initialisation process. This initialisation process
            expects to receive idles only for 6 seconds without error. If an
            error or a non-idle word is received in this time the received is
            reset and a further wait ensues. When 6 seconds of idles have been
            received the initialisation flag is set and the channel can now receive
            data. The initialisation process occurs immediately following
            configuration of the FGA. Clearly during this initialisation, and any
            others that are explicitly induced, the transmitter must be sending
            idles.

            Input parameter {\bf tx\_diff\_ctrl} specifies the differential
            output voltage and is 500 by default. Values are 400, 500, 600, 700
            or 800 which is the differential voltage in mV. The peak-to-peak
            voltage is twice this figure.

            Input parameter {\bf tx\_preemphasis} specifies the transmitter
            preemphasis and is 0 by default. values are 0 (10\%), 1 (20\%),
            2 (25\%) or 3 (33\%).

            Input parameters {\bf tx\_diff\_ctrl} and {\bf tx\_preemphasis} control
            the transmitter output waveforms to correct for different cable types
            and lengths. Any changes to the default values are usually done
            empirically based on the measured error rate.

            Input parameter {\bf loopback} is an immediate integer which specifies
            the loopback mode. If the argument is omitted or is 0, operation is
            normal, i.e. loopback is not connected. Value 1 specifies internal
            parallel mode. Value 2 specifies internal serial mode.
            
            Input parameter {\bf los\_fsm} controls the use of the MGT loss-of-sync
            finite state machine which detects when sync is held/lost. If the argument
            to this is false then {\bf rxdiag[3]} will only indicate if channel
            initialisation was successful. If {\bf los\_fsm} is true then  {\bf
            rxdiag[3]} will be asserted when the receiver has been successfully
            initialised and remain asserted while sync is OK.

            Output parameter {\bf out} is the output data queue for the receiver.
            If the receiver is not used the argument is omitted. The type can be
            struct {\bf ds16}, struct {\bf gt\_16\_e} or  struct {\bf
            gt\_16\_ee} depending on the level of error
            detail required. The {\bf k} field in each case is true if the most
            significant byte of the received word contains a K-character.
            For {\bf gt\_16\_e} or  {\bf gt\_16\_ee} the {\bf first} field
            is true then this word follows an idle, which information may be
            used to reset a block boundary. {\bf Important note - the output
            queue must not block! The destination module must read it at a rate
            that precludes it ever becoming full, if necessary reading it to
            null. Generally the read clock domain must have an equal or higher
            frequency than the user interface clock {\bf ruclk}, but this may be
            relaxed if it is known that the serial data source has a guaranteed
            maximum data rate that would allow a lower queue read clock domain.}
            For {\bf gt\_16\_e} the {\bf error} field is true if either byte
            exhibits a disparity error or not-in-table error. {\bf gt\_16\_ee}
            provides each of these errors separately for each byte.
            
            A word in which a disparity error or a not-in-table error occurs is
            written to the output queue but the error field is true. If the top
            byte of a word is a K-character the {\bf k} field is true.

            Output parameter {\bf ruclk} is the user interface clock at the base
            frequency.
            
            Output parameter {\bf rxdiag} provides five diagnostic outputs for the
            receiver. These are -
            \blist
            \item   diag[0] - rxcommadet, which is true when a positive or
                    negative comma is being received
            \item   diag[1] - rxdisperr[0] or rxnotintable[0], which is true
                    when a disparity error or a not-in-table error is
                    flagged for the current upper byte of the 16-bit output
            \item   diag[2] - rxdisperr[1] or rxnotintable[1], which is true
                    when a disparity error or a not-in-table error is
                    flagged for the current lower byte of the 16-bit output
            \item   diag[3] - is true if the receiver is synchronised, as
                    indicated by the receiver synchronisation FSM \footnote{if
                    parameter los\_fsm is true.} , and the
                    initialisation flag is high
            \item   diag[4] - is true if the receiver is receiving idles
            \blend
            These outputs might be connected to status LEDs to indicate that the
            receiver is locked and functioning correctly. {\bf diag[0]} will be true
            while the receiver is receiving idles but will be false when data are
            received.

            Output parameter {\bf txdiag} provides two diagnostic outputs for the
            transmitter. These are -
            \blist
            \item   diag[0] - is true if the transmitter is sending idles
            \item   diag[1] - is true if the transmitter is sending a word in which
                    the top byte is a control character
            \blend


    \subsection{library procedures}

        \subsubsection{mgt - create an MGT high-speed serial I/O channel}\label{sec:lpmgt}

            {\bf Library: xilinx/mgt.3pl}

            {\bf Prototype:} \\
            \vsp
            \begin{lstlisting}[gobble=12, language=threepl]
               global procedure
               mgt (
                   // outputs
                   value(rxdata),
                   value(rxchariscomma, "[]log"),
                   value(rxcharisk, "[]log"),
                   value(rxcommadet, "log"),
                   value(rxlossofsync, "uint:2"),
                   value(rxnotintable, "[]log"),
                   value(rxbufstatus, "[2]log"),
                   value(rxcheckingcrc, "log"),
                   value(rxclkcorcnt, "uint:3"),
                   value(rxcrcerr, "log"),
                   value(rxdisperr, "[]log"),
                   value(rxrealign, "log"),
                   value(rxrecclk, "log"),
                   value(rxrundisp, "[]log"),
                   value(txbuferr, "log"),
                   value(txkerr, "[]log"),
                   value(txrundisp, "[]log"),
                   value(chbondo, "[4]log"),
                   value(chbonddone, "log")
                   <-
                   // attributes
                   str(loc),
                   str(txn_pin),
                   str(txp_pin),
                   str(rxn_pin),
                   str(rxp_pin),
                   log(align_comma_msb),
                   int(chan_bond_limit),
                   str(chan_bond_mode),
                   int(chan_bond_offset),
                   log(chan_bond_one_shot),
                   int(chan_bond_seq_1_1),
                   int(chan_bond_seq_1_2),
                   int(chan_bond_seq_1_3),
                   int(chan_bond_seq_1_4),
                   int(chan_bond_seq_2_1),
                   int(chan_bond_seq_2_2),
                   int(chan_bond_seq_2_3),
                   int(chan_bond_seq_2_4),
                   log(chan_bond_seq_2_use),
                   int(chan_bond_seq_len),
                   int(chan_bond_wait),
                   log(clk_cor_insert_idle_flag),
                   log(clk_cor_keep_idle),
                   int(clk_cor_repeat_wait),
                   int(clk_cor_seq_1_1),
                   int(clk_cor_seq_1_2),
                   int(clk_cor_seq_1_3),
                   int(clk_cor_seq_1_4),
                   int(clk_cor_seq_2_1),
                   int(clk_cor_seq_2_2),
                   int(clk_cor_seq_2_3),
                   int(clk_cor_seq_2_4),
                   log(clk_cor_seq_2_use),
                   int(clk_cor_seq_len),
                   log(clk_correct_use),
                   int(comma_10b_mask),
                   int(crc_end_of_pkt),
                   str(crc_format),
                   int(crc_start_of_pkt),
                   log(dec_mcomma_detect),
                   log(dec_pcomma_detect),
                   log(dec_valid_comma_only),
                   int(mcomma_10b_value),
                   log(mcomma_detect),
                   int(pcomma_10b_value),
                   log(pcomma_detect),
                   int(ref_clk_v_sel),
                   log(rx_crc_use),
                   int(rx_data_width),
                   log(rx_decode_use),
                   int(rx_los_invalid_incr),
                   int(rx_los_threshold),
                   log(rx_loss_of_sync_fsm),
                   log(serdes_10b),
                   int(termination_imp),
                   int(tx_crc_force_value),
                   log(tx_crc_use),
                   int(tx_data_width),
                   int(tx_diff_ctrl),
                   int(tx_preemphasis),
                   // Inputs
                   value(chbondi, "[4]log"),
                   value(enchansync, "log"),
                   value(enmcommaalign, "log"),
                   value(enpcommaalign, "log"),
                   value(loopback, "uint:2"),
                   value(powerdown, "log"),
                   value(refclk, "log"),
                   value(refclk2, "log"),
                   value(brefclk, "log"),
                   value(brefclk2, "log"),
                   value(refclksel, "log"),
                   value(rxpolarity, "log"),
                   value(rxreset, "log"),
                   clock(rxusrclk),
                   clock(rxusrclk2),
                   value(txbypass8b10b, "[]log"),
                   value(txchardispmode, "[]log"),
                   value(txchardispval, "[]log"),
                   value(txcharisk, "[]log"),
                   value(txdata),
                   value(txforcecrcerr, "log"),
                   value(txinhibit, "log"),
                   value(txpolarity, "log"),
                   value(txreset, "log"),
                   clock(txusrclk),
                   clock(txusrclk2)

               ) {}
            \end{lstlisting}

            {\bf Output parameters:} \\
            See Xilinx documentation.

            {\bf Input parameters:} \\
            See Xilinx documentation.

            {\bf In-line target execution:} \\
            no

            {\bf description:} \\
            This procedure creates an MGT high speed serial I/O channel.
            For a simpler interface
            see {\bf gt\_custom\_16} \autorefph{sec:lmgtcustom16} above.



    \subsection{library functions}
    
        None.


    \subsection{library classes}
    
        None.
